% Chapter 04: Lowering
% The Kit compiler: transforming domain syntax into core primitives

\chapter{Lowering}
\label{ch:lowering}

% ============================================================
\section{The Lowering Contract}
\label{sec:lowering-contract}
% ============================================================

\textbf{Lowering} is the process of compiling Kit-authored runbook YAML into
core \texttt{runbook/v1} YAML. The Kit compiler is a pure function:

\begin{verbatim}
Kit YAML (stdin) → Kit Compiler → Core YAML (stdout)
\end{verbatim}

\paragraph{Input:} A runbook with \texttt{apiVersion: runbook/v1+<kit-suffix>}
that uses Kit-specific step types (e.g., \texttt{approval-request},
\texttt{evidence-capture}).

\paragraph{Output:} A valid \texttt{runbook/v1} document that uses only core
step types (\texttt{cli}, \texttt{tool}, \texttt{manual}, \texttt{branch},
\texttt{iterate}) and core fields.

\paragraph{Key properties:}

\begin{itemize}
  \item \textbf{Purity:} The compiler has no side effects. It does not invoke
    binaries, call APIs, write to disk (except stdout), or access the network.
  \item \textbf{Determinism:} The same input always produces the same output.
    This is critical for testing and reproducibility.
  \item \textbf{Correctness:} The lowered runbook must be semantically
    equivalent to the Kit-authored runbook. Lowering can add steps
    (synthesized steps) but must preserve the author's intent.
\end{itemize}

% ============================================================
\section{Purity Requirement}
\label{sec:purity}
% ============================================================

The Kit compiler MUST be a pure transformation. This means:

\begin{itemize}
  \item No filesystem reads (except stdin).
  \item No filesystem writes (except stdout).
  \item No network requests.
  \item No subprocess invocations.
  \item No environment variable dependencies (except for optional debug flags).
  \item No non-deterministic operations (random numbers, timestamps, etc.).
\end{itemize}

\paragraph{Why purity matters:}

\begin{enumerate}
  \item \textbf{Testability:} Pure functions are easy to test. You provide
    input, assert output, done.
  \item \textbf{Reproducibility:} The same runbook lowered today and lowered
    tomorrow produces identical output. This is essential for version control
    and CI systems.
  \item \textbf{Auditability:} Lowering is transparent. You can inspect the
    lowered output and trace it back to the Kit input line-by-line.
  \item \textbf{Portability:} A pure compiler can run anywhere (local laptop,
    CI runner, airgapped environment) without external dependencies.
\end{enumerate}

If your Kit needs runtime state (e.g., fetching user lists from an API),
that belongs in a \textbf{Tool} or \textbf{Extension}, not in the compiler.

% ============================================================
\section{What Lowering Must Preserve}
\label{sec:must-preserve}
% ============================================================

The lowering MUST preserve the following:

\begin{description}
  \item[Step order] --- Steps execute in the same order as authored (unless the
    author explicitly uses \texttt{branch} or \texttt{iterate}).
  \item[Variable bindings] --- If the Kit-authored runbook sets a variable
    \texttt{X}, the lowered runbook must set \texttt{X} with the same value.
  \item[Governance hooks] --- If a Kit step specifies governance fields (e.g.,
    \texttt{approval}, \texttt{allow}, \texttt{redact}), those must be
    preserved in the lowered core steps.
  \item[Evidence slots] --- If a Kit step captures evidence, the lowered core
    step must have the corresponding evidence field.
  \item[Step metadata] --- The lowered step should preserve \texttt{id},
    \texttt{description}, and any other metadata from the Kit step.
\end{description}

% ============================================================
\section{What Lowering May Add}
\label{sec:may-add}
% ============================================================

The lowering MAY add the following:

\begin{description}
  \item[Synthesized steps] --- Steps that were not authored but are necessary
    for correctness. Example: a Kit might synthesize a setup step before the
    first user-authored step.
  \item[Variable injections] --- The compiler may inject variables into the
    runbook's \texttt{vars} section (e.g., Kit version, lowering timestamp).
  \item[Pre/post hooks] --- The compiler may wrap user steps with pre-step or
    post-step hooks (implemented as additional core steps).
  \item[Governance field population] --- If a Kit step implies governance (e.g.,
    \texttt{approval-request} implies \texttt{approval: manual}), the compiler
    populates those fields in the lowered output.
\end{description}

\paragraph{Synthesized step ID conventions:}

If your compiler synthesizes steps, use a naming convention that makes it clear
they are synthetic:

\begin{itemize}
  \item Prefix with \texttt{\_\_kit\_\_} (e.g., \texttt{\_\_kit\_\_setup}).
  \item Or suffix with \texttt{\_\_synthetic} (e.g., \texttt{init\_\_synthetic}).
\end{itemize}

This helps debuggers and trace viewers distinguish authored steps from
synthesized steps.

% ============================================================
\section{The Compiler Binary Interface}
\label{sec:compiler-interface}
% ============================================================

\paragraph{Invocation:}

\begin{verbatim}
<kit-compiler-binary> compile < input.yaml > output.yaml
echo $?   # Exit code: 0 = success, 1 = error
\end{verbatim}

\paragraph{Stdin:} Kit-authored runbook YAML.

\paragraph{Stdout:} Lowered \texttt{runbook/v1} YAML (if successful) OR error
messages in JSON format (if failed).

\paragraph{Stderr:} Diagnostic messages, debug logs (optional, not parsed by gert).

\paragraph{Exit codes:}

\begin{description}
  \item[\texttt{0}] --- Lowering succeeded. Stdout contains valid
    \texttt{runbook/v1} YAML.
  \item[\texttt{1}] --- Lowering failed. Stdout contains a JSON array of error
    messages (same format as validator output, \S\ref{sec:error-format}).
  \item[\texttt{2+}] --- Internal compiler error (panic, unexpected failure).
    Stderr may contain a stack trace.
\end{description}

\paragraph{Error output format:}

If lowering fails (exit code 1), write a JSON array of error messages to stdout:

\begin{minted}{json}
[
  {
    "file": "runbook.yaml",
    "line": 12,
    "severity": "error",
    "message": "approval-request step requires at least one approver",
    "rule": "compliance/lowering"
  }
]
\end{minted}

% ============================================================
\section{Example: Lowering \texttt{approval-request}}
\label{sec:example-lowering}
% ============================================================

Let's lower a Kit-authored \texttt{approval-request} step into a core
\texttt{manual} step with governance fields populated.

\paragraph{Kit input:}

\begin{minted}{yaml}
apiVersion: runbook/v1+compliance/v1
meta:
  title: "Deploy to production"

steps:
  - id: approve-deploy
    type: approval-request
    description: "Get approval from ops team"
    approvers:
      - alice@acme.example
      - bob@acme.example
    timeout: 2h
    evidence-required: true
\end{minted}

\paragraph{Lowered output:}

\begin{minted}{yaml}
apiVersion: runbook/v1
meta:
  title: "Deploy to production"
  x-kit:lowered-from:
    kit: gert.compliance
    version: "1.2.0"

steps:
  - id: approve-deploy
    type: manual
    description: "Get approval from ops team"
    prompt: |
      Approval required from:
        - alice@acme.example
        - bob@acme.example
      
      Timeout: 2h
      Evidence required: Yes
    
    approval:
      required: true
      approvers:
        - alice@acme.example
        - bob@acme.example
      timeout: 7200  # 2h in seconds
    
    evidence:
      required: true
      type: text
      prompt: "Provide evidence for this approval (required)"
    
    on-timeout:
      action: reject
      message: "Approval timed out after 2h"
\end{minted}

\paragraph{Transformation notes:}

\begin{itemize}
  \item The Kit step type \texttt{approval-request} is lowered to core type
    \texttt{manual}.
  \item Kit fields (\texttt{approvers}, \texttt{timeout},
    \texttt{evidence-required}) are mapped to core governance fields
    (\texttt{approval}, \texttt{evidence}, \texttt{on-timeout}).
  \item The compiler synthesized a \texttt{prompt} field to guide the operator.
  \item Metadata was added to \texttt{meta.x-kit:lowered-from} to record which
    Kit and version performed the lowering.
\end{itemize}

% ============================================================
\section{Example: Compiler Pseudocode}
\label{sec:compiler-pseudocode}
% ============================================================

Here is pseudocode for a minimal Kit compiler (Go):

\begin{minted}{go}
package main

import (
  "encoding/json"
  "fmt"
  "gopkg.in/yaml.v3"
  "os"
)

func main() {
  var runbook map[string]interface{}
  if err := yaml.NewDecoder(os.Stdin).Decode(&runbook); err != nil {
    reportError("Failed to parse input YAML", 1)
    os.Exit(1)
  }

  lowered, err := lowerRunbook(runbook)
  if err != nil {
    reportError(err.Error(), 1)
    os.Exit(1)
  }

  yaml.NewEncoder(os.Stdout).Encode(lowered)
}

func lowerRunbook(rb map[string]interface{}) (map[string]interface{}, error) {
  // Change apiVersion from runbook/v1+compliance/v1 to runbook/v1
  rb["apiVersion"] = "runbook/v1"

  // Add lowering metadata
  meta := rb["meta"].(map[string]interface{})
  meta["x-kit:lowered-from"] = map[string]string{
    "kit":     "gert.compliance",
    "version": "1.2.0",
  }

  // Lower each step
  steps := rb["steps"].([]interface{})
  for i, step := range steps {
    s := step.(map[string]interface{})
    loweredStep, err := lowerStep(s)
    if err != nil {
      return nil, fmt.Errorf("step %d: %w", i, err)
    }
    steps[i] = loweredStep
  }

  return rb, nil
}

func lowerStep(step map[string]interface{}) (map[string]interface{}, error) {
  stepType := step["type"].(string)

  switch stepType {
  case "approval-request":
    return lowerApprovalRequest(step)
  case "evidence-capture":
    return lowerEvidenceCapture(step)
  default:
    // Core step types pass through unchanged
    return step, nil
  }
}

func lowerApprovalRequest(step map[string]interface{}) (map[string]interface{}, error) {
  approvers := step["approvers"].([]interface{})
  if len(approvers) == 0 {
    return nil, fmt.Errorf("approval-request requires at least one approver")
  }

  timeout := step["timeout"].(string)  // e.g., "2h"
  timeoutSec := parseTimeout(timeout)  // Convert to seconds

  lowered := map[string]interface{}{
    "id":          step["id"],
    "type":        "manual",
    "description": step["description"],
    "prompt":      buildApprovalPrompt(step),
    "approval": map[string]interface{}{
      "required":  true,
      "approvers": approvers,
      "timeout":   timeoutSec,
    },
  }

  if step["evidence-required"] == true {
    lowered["evidence"] = map[string]interface{}{
      "required": true,
      "type":     "text",
      "prompt":   "Provide evidence for this approval (required)",
    }
  }

  return lowered, nil
}

func buildApprovalPrompt(step map[string]interface{}) string {
  approvers := step["approvers"].([]interface{})
  timeout := step["timeout"].(string)
  var prompt string
  prompt += "Approval required from:\n"
  for _, a := range approvers {
    prompt += fmt.Sprintf("  - %s\n", a)
  }
  prompt += fmt.Sprintf("\nTimeout: %s\n", timeout)
  if step["evidence-required"] == true {
    prompt += "Evidence required: Yes\n"
  }
  return prompt
}

func parseTimeout(t string) int {
  // Parse "2h" → 7200 seconds (simplified; real impl would handle all units)
  // ...
  return 7200
}

func reportError(msg string, line int) {
  errors := []map[string]interface{}{
    {
      "file":     "<stdin>",
      "line":     line,
      "severity": "error",
      "message":  msg,
      "rule":     "compliance/lowering",
    },
  }
  json.NewEncoder(os.Stdout).Encode(errors)
}
\end{minted}

% ============================================================
\section{Testing Your Compiler}
\label{sec:testing-compiler}
% ============================================================

The compiler MUST be tested with fixture runbooks and expected outputs. See
Chapter~\ref{ch:testing-contracts} for the full testing workflow. At a minimum:

\begin{enumerate}
  \item Create a fixture runbook: \texttt{test/fixtures/basic-approval.yaml}
  \item Run your compiler: \texttt{cat basic-approval.yaml | ./bin/gert-kit-compliance compile > output.yaml}
  \item Assert that \texttt{output.yaml} matches the expected lowered form.
  \item Validate the lowered output against the core \texttt{runbook/v1} schema.
\end{enumerate}

Automate this with a test script or a \texttt{Makefile} target.
